% =============================================================================
% ARTICLE RECORD — faithful LaTeX transcription of a published journal article
% -----------------------------------------------------------------------------
% bibkey:    <AuthorYear e.g. FalmagneLundberg2000>
% authors:   <Full Name; Full Name>
% title:     <exact title>
% journal:   <journal name>
% volume:    <vol>
% year:      <year>
% pages:     <ppp--ppp>
% doi:       <doi or empty>
% issn:      <issn or empty>
% msc:       <subject classification codes, if any>
% keywords:  <k1; k2; k3>
% source_pdf: <year_author_title.pdf>
%
% PROVENANCE
%   Transcribed from the <born-digital text layer | scanned pages via OCR> of the
%   source PDF (N pp.). Method: <read pages as images -> transcribe LaTeX -> pdflatex
%   compile (0 errors) -> visual diff of rendered output vs source pages>. Numbered
%   equations carry the article's OWN numbers via \tag{}. Connecting proof prose is
%   condensed; every numbered display equation is transcribed. Original
%   inconsistencies are FLAGGED in the Editorial notes section, never silently fixed.
% =============================================================================
\documentclass[11pt]{article}
\usepackage[margin=1in]{geometry}
\usepackage{amsmath,amssymb,amsthm}
\usepackage{array}
% \DeclareMathOperator{\Dom}{Dom}   % declare any operators the paper uses
\newenvironment{stmt}{\itshape}{\par}   % for theorem/definition statements
\begin{document}

\begin{center}
{\large\bfseries <TITLE>}\\[4pt]
<Authors>\\[2pt]
{\itshape <Journal>} \textbf{<vol>} (<year>) <pages>
\end{center}

\noindent\textbf{Summary (verbatim).} <abstract, quoted verbatim>

\bigskip\hrule\bigskip

\section*{Core definitions / equations (Introduction)}

\begin{equation}
% <equation> \tag{1}
\end{equation}

\section*{1.\ <section title>}

\paragraph{1.1.\ Theorem.}
\begin{stmt}
<theorem statement, verbatim in content>
\end{stmt}
\begin{equation}
% <equation> \tag{6}
\end{equation}

\paragraph{Proof (condensed).}
<key numbered equations of the proof, prose bridges summarised>
\hfill$\square$

\section*{2.\ <main result section>}

\paragraph{2.1.\ Definition.}
\begin{stmt}<...>\end{stmt}

\paragraph{2.N.\ Theorem (main result).}
\begin{stmt}<...>\end{stmt}

\section*{References (verbatim)}
\begin{enumerate}
\item[{[1]}] <reference, verbatim>
\end{enumerate}

\bigskip\hrule\bigskip

\section*{Editorial / transcription notes (not part of the original paper)}
\begin{enumerate}
\item \textbf{<flagged issue>.} <e.g. original theorem-number typo; ambiguous glyph;
what was condensed; anything a reader should know. Include your confidence.>
\end{enumerate}

\end{document}
